% ------------------------------------------------------------------------
% file `figures_planes-exercise-7-solution-body.tex'
%   in folder `xsim/'
%
%     solution of type `exercise' with id `7'
%
% generated by the `solution' environment of the
%   `xsim' package v0.21 (2022/02/12)
% from source `figures_planes' on 2026/08/04 on line 428
% ------------------------------------------------------------------------
\begin{enumerate}
    \item Un angle dont l'amplitude est comprise entre $0^\circ$ et $90^\circ$ est un \textbf{angle aigu}. L'angle de $50^\circ$ est donc un angle aigu.
    \item L'élève doit choisir la graduation $50^\circ$. La graduation $130^\circ$ correspondrait à un angle obtus (plus grand qu'un angle droit de $90^\circ$). Pour ne pas se tromper, l'élève doit toujours vérifier visuellement si son angle est "fermé" (aigu, $<90^\circ$) ou "ouvert" (obtus, $>90^\circ$).
\end{enumerate}
